\documentclass[11pt]{article}
\input{../../shared/project_style.tex}
\rhead{Basics 25 Textbook Draft}

\begin{document}
\DocTitle{Basics 25: Conservation Laws in Magnetohydrodynamics}{IV - Magnetohydrodynamics}{Integral balances of mass, momentum, energy, magnetic flux, and helicity}

\begin{overviewbox}
\textbf{Textbook draft, version 1.0.} Conservation laws are both physical principles and numerical acceptance tests. This chapter develops the local and integral balances of ideal MHD, distinguishes exact invariants from boundary-dependent ones, and shows how conservation constrains multiphysics coupling. This chapter is designed for a reader with no prior plasma-physics training. It distinguishes exact identities from physical approximations and connects the central equations to later simulation modules.
\end{overviewbox}

\ProjectTOC

\section{Learning objectives}
After completing this note, the reader should be able to:
\begin{itemize}[leftmargin=1.6em]
\item Convert local conservation equations into integral balances.
\item Interpret momentum flux and Maxwell stress.
\item Derive ideal-MHD total energy conservation.
\item State the conditions for magnetic-flux and helicity conservation.
\item Distinguish global conservation from local redistribution.
\item Design conservation diagnostics for coupled solvers.
\end{itemize}

\section{Prerequisites and reading strategy}
\begin{itemize}[leftmargin=1.6em]
\item Basics 3, 12, 14, and 24.
\end{itemize}
On a first reading, emphasize physical meaning and dimensional checks. On a second reading, reproduce the derivations. Attempt the computational laboratory only after the limiting cases are understood.

\section{Generic local and integral balance}

A local conservation law is
\begin{equation}
\partial_tU+\nabla\cdot\mathbf F=S.
\end{equation}
Integrating over a fixed control volume $V$ gives
\begin{equation}
\frac{d}{dt}\int_VU\,dV
=-\oint_{\partial V}\mathbf F\cdot d\mathbf A+\int_VS\,dV.
\end{equation}
A quantity can change inside $V$ because it crosses the boundary or because a true source acts. Numerical diffusion should not be silently reported as a physical source.


\section{Mass conservation}

Ideal MHD satisfies
\begin{equation}
\partial_t\rho+\nabla\cdot(\rho\mathbf u)=0.
\end{equation}
For an impermeable or periodic boundary,
\begin{equation}
\frac{d}{dt}\int_V\rho\,dV=0.
\end{equation}
In a multiphysics code, particle birth, wall loss, or source modules must appear explicitly so that the same particles are not removed or deposited twice.


\section{Momentum and stress transport}

The momentum equation can be written
\begin{equation}
\partial_t(\rho\mathbf u)+\nabla\cdot\mathsf T=0,
\end{equation}
where
\begin{equation}
\mathsf T=
\rho\mathbf u\mathbf u+
\left(p+\frac{B^2}{2\mu_0}\right)\mathsf I
-\frac{\mathbf B\mathbf B}{\mu_0}.
\end{equation}
The tensor contains advective momentum flux, gas pressure, magnetic pressure, and directional magnetic tension. Total momentum is conserved only when boundary tractions and external forces vanish.


\section{Total energy}

The conserved ideal-MHD energy density is
\begin{equation}
E=\frac{\rho u^2}{2}+\frac{p}{\gamma-1}+\frac{B^2}{2\mu_0}.
\end{equation}
Its flux is
\begin{equation}
\mathbf F_E=
\left(E+p+\frac{B^2}{2\mu_0}\right)\mathbf u
-\frac{\mathbf B(\mathbf u\cdot\mathbf B)}{\mu_0}.
\end{equation}
The magnetic contribution can also be interpreted through Poynting flux $\mathbf E\times\mathbf B/\mu_0$ with ideal $\mathbf E=-\mathbf u\times\mathbf B$. Energy may exchange among kinetic, internal, and magnetic reservoirs while total energy remains constant.


\section{Magnetic flux conservation}

For a material surface $S(t)$ moving with ideal-MHD velocity,
\begin{equation}
\frac{d}{dt}\int_{S(t)}\mathbf B\cdot d\mathbf A=0.
\end{equation}
This is Alfvén's flux-freezing theorem. It constrains magnetic topology under smooth ideal evolution. It does not imply that $B$ is constant: compression and stretching alter field strength while preserving flux through the moving surface.


\section{Magnetic and cross helicity}

Magnetic helicity is
\begin{equation}
H_M=\int_V\mathbf A\cdot\mathbf B\,dV.
\end{equation}
Its gauge-invariant interpretation requires suitable closed or relative-helicity boundary conditions. Cross helicity,
\begin{equation}
H_C=\int_V\mathbf u\cdot\mathbf B\,dV,
\end{equation}
is conserved in restricted ideal barotropic settings. These invariants characterize topology and alignment, but their conservation conditions are narrower than those of mass.


\section{Entropy and irreversible physics}

For adiabatic ideal flow,
\begin{equation}
\frac{D}{Dt}\left(\frac{p}{\rho^\gamma}\right)=0.
\end{equation}
Resistivity, viscosity, shocks, heat conduction, and collisions generate entropy. A mathematically conservative shock-capturing method can conserve total energy while converting resolved kinetic or magnetic energy into internal energy.


\section{Conservation across module interfaces}

Suppose a fast-particle module reports lost energy $\Delta E_{\rm fast}<0$. A wall or thermal module should receive exactly $-\Delta E_{\rm fast}$ minus any explicitly exported energy:
\begin{equation}
\Delta E_{\rm fast}+\Delta E_{\rm wall}+\Delta E_{\rm other}=0.
\end{equation}
Similar ledgers apply to particle number and momentum. Conservative remapping and synchronized timestamps are therefore part of physics correctness, not merely software design.


\section{Computational laboratory}
Run the reduced Alfvén benchmark or a simple periodic MHD wave and track kinetic, magnetic, internal, and total energy. Perturb interpolation between two grids and measure the conservation defect. Implement a module-level particle and energy ledger.

\section{Summary}
\begin{itemize}[leftmargin=1.6em]
\item Conservation laws separate transport through boundaries from true sources.
\item MHD momentum is transported by fluid and Maxwell stresses.
\item Total energy can remain conserved while its physical reservoirs exchange energy.
\item Flux and helicity invariants require ideal dynamics and appropriate boundaries.
\item Coupled simulation modules require explicit particle, momentum, and energy ledgers.
\end{itemize}

\SymbolsAcronymsSection
\begin{symtable}
$U$ & conserved density & Quantity per volume. \\
$\mathbf F$ & flux & Transport rate per area. \\
$\mathsf T$ & momentum-flux tensor & Fluid plus magnetic stress. \\
$H_M$ & magnetic helicity & $\int\mathbf A\cdot\mathbf B\,dV$. \\
$H_C$ & cross helicity & $\int\mathbf u\cdot\mathbf B\,dV$. \\
$S$ & source density & Nonconservative production or removal. \\
\end{symtable}

\section{Exercises}
\begin{enumerate}[leftmargin=1.8em]
\item Derive the integral mass balance from the local equation.
\item Identify every term in the MHD momentum-flux tensor.
\item Rewrite the electromagnetic part of the energy flux using ideal Ohm's law.
\item State boundary conditions required for global momentum and energy conservation.
\item Design a conservation test for alpha-particle energy transferred to thermal plasma and wall loads.
\end{enumerate}

\section{References}
\begin{thebibliography}{99}
\bibitem{ref1} G. B. Arfken, H. J. Weber, and F. E. Harris, \emph{Mathematical Methods for Physicists}, 7th ed., Academic Press (2013).
\bibitem{ref2} G. Strang, \emph{Linear Algebra and Its Applications}, 4th ed., Brooks/Cole (2006).
\bibitem{ref3} W. A. Strauss, \emph{Partial Differential Equations: An Introduction}, 2nd ed., Wiley (2007).
\bibitem{ref4} D. J. Griffiths, \emph{Introduction to Electrodynamics}, 4th ed., Cambridge University Press (2017).
\bibitem{ref5} J. D. Jackson, \emph{Classical Electrodynamics}, 3rd ed., Wiley (1998).
\bibitem{ref6} J. P. Freidberg, \emph{Ideal MHD}, Cambridge University Press (2014).
\bibitem{ref7} J. P. Goedbloed, R. Keppens, and S. Poedts, \emph{Advanced Magnetohydrodynamics}, Cambridge University Press (2010).
\bibitem{ref8} H. Goedbloed and S. Poedts, \emph{Principles of Magnetohydrodynamics}, Cambridge University Press (2004).
\end{thebibliography}

\end{document}
